\documentclass{article}
\title{Practice Problem 2.17}
\author{CSAPP}
\usepackage{tabularx}
\usepackage{multirow}
\usepackage{physics}
\usepackage{array}
\usepackage[table]{xcolor}
\usepackage[margin=1in]{geometry}
\usepackage{booktabs} % For better table rules

\begin{document}

\maketitle

Assuming $w = 4$, we can assign a numeric value to each possible hexadecimal digit, assuming either an unsigned or a two's-complement interpretation. Fill in the following table according to these interpretations by writing out the nonzero powers of $2$.

\begin{table}[h]
\centering
\begin{tabular}{|l|l|l|l|}
    \hline
    \multicolumn{4}{|c|}{$\vec{x}$} \\
    \hline
    \textbf{Hexadecimal} & \textbf{Binary} & $B2U_4(\vec{x})$ & $B2T_4(\vec{x})$ \\
    \hline
    \texttt{0xA} & \texttt{1010} & $2^3 + 2^1 = 10$ & $-2^3 + 2^1 = -6$ \\
    \hline
    \texttt{0x1} & \texttt{0001} & $2^0 = 1$ & $2^0 = 1$ \\
    \hline
    \texttt{0xB} & \texttt{1011} & $2^3 + 2^1 + 2^0 = 11$ & $-2^3 + 2^1 + 2^0 = -5$ \\
    \hline
    \texttt{0x2} & \texttt{0010} & $2^1 = 2$ & $2^1 = 2$ \\
    \hline
    \texttt{0x7} & \texttt{0111} & $2^2 + 2^1 + 2^0 = 7$ & $2^2 + 2^1 + 2^0 = 7$ \\
    \hline
    \texttt{0xC} & \texttt{1100} & $2^3 + 2^2 = 12$ & $-2^3 + 2^2 = -4$ \\
    \hline
\end{tabular}
\end{table}

\end{document}